\section{The three-letter scalar lift with two up positions}

% BEGIN A83U1
\paragraph{A83U1. The denominator of the unique down residue.}
On the mixed orbit of U83M1 take a length-three arc beginning up.
Its signs from $c-2$ to $c$ are up, down, up. Write
\[
 a=s_{c-2}\le w,\quad b=s_c,\quad D=s_{c-1}+S,\quad
 Q=P_0=v+p,\quad Y=P_1,\quad Z=P_2,\quad A=O[[v]].
\]
Both later sizes are in $[B,B+\eta]$, and
$20w,20\eta,20S\le B$ on the late family. For a native scalar
$f$, $N=Q^bZ^a f$ is integral and there is a constant $z\in L$
such that $N-zQ^b Z^a\in Y^D A[1/p]$. Put
$s=\max(0,\lceil-v_p(z)\rceil)$, taking $s=0$ at $z=0$.
Then
\[
                  s\le \max(0,a-D/p)+4U .
\]

For $v_p(z)<0$, write $n=\lfloor-v_p(z)\rfloor$ and
$t=\lceil b/(p-1)\rceil$. In the finite free $O$-module
$A/Y^D$, the class of $z Q^b Z^a$ is integral, so $Q^b Z^a$
belongs to $(p^n,Y^D)$. As $Y-p$ is a unit times $vQ^{p-1}$,
the multiplier $(Y-p)^t/Q^b$ is integral. The ring $A$ is
finite free over $O[[Y]]$, with basis containing $1$, by
preparation of $P_1(v)-Y$ of degree $p$. Moreover $Z$ is a
series in $Y$, of the form $p\varepsilon(Y)+Y^p U_1(Y)$ with
both coefficients units, by Frobenius applied to
$P_1=p+vQ^{p-1}U_0$. Thus
$(Y-p)^t Z^a\in(p^n,Y^D)$ already in $O[[Y]]$.
Evaluation at $Y=0$ gives $n\le t+a$. Expansion of
$Y^{t+n}=((Y-p)+p)^{t+n}$ shows that it belongs to
$((Y-p)^t,p^n)$, and hence
\[
              Z^a\in(p^n,Y^{D-t-n})
\]
by cancellation of $Y^{t+n}$ modulo $p^n$. The exponent is
positive for all sufficiently late indices, uniformly:
$t+n\le a+2t\le B/20+2(B+\eta)/(p-1)+2<B\le D$.
For $n=0$ the desired bound is immediate.

Put $j=\min(a,\lfloor(D-t-n-1)/p\rfloor)$. In the binomial
expansion of $(p\varepsilon+Y^p U_1)^a$ only the terms indexed
$0\le i\le j$ survive modulo $Y^{D-t-n}$. Their scalar
coefficients are $p^{a-i}\binom ai$. Each belongs to $p^nO$:
otherwise divide the sum by a constant of the smallest coefficient
valuation and reduce modulo the uniformizer. The unit terms have
distinct first degrees $pi<D-t-n$, so the earliest survives.
In particular
$n\le a-j+v_p\binom aj$. If $j=a$, then $n=0$. Otherwise
$v_p\binom aj\le\log_p(2h)=L_0$ by the factorial formula, so
\[
 s\le n+1\le a-D/p+(a+2t+1)/p+2+L_0 .
\]
Here $a/p\le w/p$, $2t/p\le(11/5)h/p^2+2/p$ for $p\ge11$.
The claimed $4U$ allowance follows for the original, possibly
ramified, coefficient DVR.
% END A83U1

% BEGIN A83U2
\paragraph{A83U2. Prescribing a constant at the unused up position.}
Keep the same $a,D,Q,Y,Z$. There are $A_1,B_1\in A$ with
\[
 Z^a=p^d A_1+Y^D Q^\eta B_1,\qquad
 d=\max(0,\lfloor a-D/p-\eta/p^2\rfloor).
\]
Indeed $Z=p\varepsilon+Y^pU_1$ and $Y=p\varepsilon_0+Q^p U_0'$
with all four series integral. The $j$th summand of the first
expansion has the factor $p^{a-j}Y^{pj}$. If $pj<D$, this
already has at least $a-D/p$ factors of $p$. Otherwise extract
$Y^D$ and expand $Y^{pj-D}$ by the second formula. Terms
surviving modulo $Q^\eta$ use fewer than $\eta/p$ factors of
$Q^p$. Consequently the remaining power of $p$ is at least
\[
 a-j+\max(0,pj-D-\eta/p)\ \ge\ a-D/p-\eta/p^2 .
\]
The inequality follows on either side of $pj=D+\eta/p$.
This also covers $\eta=0$, when the extracted terms all vanish
in the indicated ideal.

Put $k=\max(0,s-d)$, so $k\le4U+1+\eta/p^2\le6U$ by A83U1.
With $x=p^s z\in O$, form
\[
      f_a=p^{k+d-s}x A_1/Z^a,\qquad f_b=p^kf-f_a .
\]
The first numerator is integral. Over $L$, $f_a-p^k z$ vanishes
to order $D$ at every $Y$ branch and order $\eta$ at every
$Q$ branch, since $Z$ is a Taylor unit at both. The numerator
$Q^bZ^a f_b$ is integral and divisible by $Y^D$ by saturation
of distinguished division. Thus $f_b$ has the two allowable
up poles and the full down zero, and the first construction of
S83L1 lifts it at exponent $R+\sigma$.
% END A83U2

% BEGIN A83U3
\paragraph{A83U3. The terminal-up length-three comparison on every tested map.}
Write $I_0=\lceil s_I/(p-1)\rceil$, where $s_I$ is the least
Wach power for the tested surjection at the first letter $c-2$.
Fix any integral target coordinate $E_2$ satisfying the up
filtrations there and at $c$. If $a>I_0+\eta$, apply the
one-step extraction of S83L3 to the prefix ending at $c-2$.
More explicitly the target test is $E_2/Q^a$ with prefix height
$G$. The direct comparison at cost $\sigma$, lifting after
$v^{I_0}$, and the inverse comparison at cost $\sigma$ send
the source numerator to $p^{2\sigma}Q^{G-a}(Q-p)^{I_0}E_2$
modulo $Q^{2G}$. Its coefficient of degree
$G-a+I_0<G-\eta$ gives an integral predecessor vector $E_1$
in the transported $\mathrm{Fil}^a$ at this first letter, with
$F(0)E_1=p^{2\sigma}E_2$. This computation uses the prefix
starting at the original predecessor $c-3$ and hence gives the
same constant map as the full three-step frame.

Now $(\gamma-1)(f_a E_1)/v\in\mathcal N_1$. At $Z$ its pole
has order at most $a$ in the filtered direction just found.
At $Y$ and $Q$ the constant $p^k z$ is fixed at every branch,
so the difference vanishes to order at least $D$ and $\eta$,
respectively, by A83U2. These orders meet all the source tests,
since the terminal up letter has no jump below $-\eta$.
The numerator of the difference vanishes at $v=0$ and divides
integrally by $v$, a Taylor unit at the three primitive levels.
U83M6 and the functorial frame congruence of S83L1 give a source
witness for $p^{3\sigma} f_aE_2$ with target correction zero.
Add the witness for $f_b$ at a common exponent. This proves
$\mathrm{LiftsAt}$ for $f$ at $k+\max(3\sigma,R+\sigma)\le15U$.

If $a\le I_0+\eta$, decompose $p^sf=p^sz+(p^sf-p^sz)$.
The second term has integral numerator with the full down zero;
the first is an integral constant. S83L1 lifts both at
$R+\sigma$, costing $s+R+\sigma$ for $f$. By R82W1,
$s_I\le h+\eta$, and $h-D\le\eta$. Hence
\[
 a-D/p\le\eta+1+h/[p(p-1)]+\eta/(p-1)+\eta/p
                 \le\eta+3U .
\]
If $a\le D/p$ then $s\le4U$. If $\eta\le8U$ the resulting
exponent is at most $21U$. In the remaining case $a>D/p$ and
$\eta>8U$, the first, nonthick, letter of size $a$ is adjacent
to the thick middle letter, and
\[
 0<pa-B\le p(\eta+3U)+D-B\le2p\eta+S+\eta .
\]
This is within the original band of coefficient $320$ on
$p\eta+pS+B/p+a+pL_0$, so its indicator $I$ is one.
The cost in this case is at most $13U+\eta$.
In all cases the witness may be multiplied by an integral
power of $p$ to realize exactly $\lfloor40U+\eta I\rfloor$ on
the given cut. This proves the native scalar comparison for
every up, down, up arc on the late mixed family.
% END A83U3
